% Shared preamble of the HSEMHS mode documents (docs/physics/*.tex, docs/technical/*.tex).
\documentclass[10pt]{extarticle}
\usepackage[a4paper,margin=1.7cm]{geometry}
\usepackage[T1]{fontenc}\usepackage{lmodern}\usepackage{tabularx}
\usepackage{amsmath,amssymb,bm,enumitem,titlesec,xcolor,booktabs,array,tikz,graphicx,longtable}
\usepackage[hidelinks]{hyperref}
\titleformat{\section}{\large\bfseries\color{blue!50!black}}{\thesection.}{0.5em}{}
\titlespacing*{\section}{0pt}{9pt}{3pt}
\titleformat{\subsection}{\normalsize\bfseries}{\thesubsection}{0.5em}{}
\titlespacing*{\subsection}{0pt}{6pt}{2pt}
\setlist{nosep,leftmargin=1.3em}
\setlength{\parindent}{0pt}\setlength{\parskip}{4pt}
\renewcommand{\arraystretch}{1.15}
\newcommand{\B}{\mathbf B}\newcommand{\J}{\mathbf J}\newcommand{\F}{\mathbf F}\newcommand{\R}{\mathbf R}
\newcommand{\zh}{\hat{\mathbf z}}\newcommand{\Bo}{\mathbf B_{\rm obs}}\newcommand{\Lv}{\mathbf L}
\newcommand{\code}[1]{\texttt{#1}}
\newcolumntype{L}[1]{>{\raggedright\arraybackslash}p{#1}}
\newcolumntype{Y}{>{\raggedright\arraybackslash}X}
\newcommand{\doctitle}[2]{{\Large\bfseries #1}\\[2pt]{\small #2}\par\vspace{2pt}}
\newif\ifadjointcolumn\newif\ifadjointpoisson   % which part of common/adjoint.tex a document shows
\pagestyle{plain}
\AtBeginDocument{\sloppy}

\begin{document}
\doctitle{mhs-poisson-logm: magnetohydrostatic correction on a column-mass grid}{Physics and method of
\code{run.mode = "mhs-poisson-logm"} in HSEMHS, 25 Nov 2026. The solver of \code{mhs-poisson-ltau} on a mass axis: read that document
for the derivations; this one gives the differences.}

\begin{abstract}
\noindent With $T$ and $\B$ given on surfaces of constant column mass $m$, this mode computes the gas pressure on every surface and
the height of every surface in every column so that pressure, gravity and the Lorentz force of the given field balance as well as
possible in all three directions, below a cut in $m$; above the cut the columns keep the hydrostatic pressure $P=gm$. It is the
least-squares problem of \code{mhs-poisson-ltau} with the optical depth replaced by the column mass.
\end{abstract}

\section{From the optical-depth axis to the mass axis}
The mass above a point changes as $dm=-\rho\,dz$, the optical depth as $d\tau=-\chi\,dz$. Every equation of \code{mhs-poisson-ltau}
that involves the axis therefore holds with $\tau\to m$ and $\chi\to\rho$:
\begin{itemize}
\item \textbf{Heights:} $z_{k+1}=z_k-\tfrac12\Delta\ln m_k\,(a_k+a_{k+1})$ with $a=\mu m/\rho$ (on the $\tau$ axis $a=\mu\tau/\chi$).
\item \textbf{Vertical balance along the axis:} with the Lorentz force, $dP/dz=-\rho g+L_z$ becomes
$d\ln P/d\ln m=\mu m\,(g-L_z/\rho)/P$; without a field the solution is $P=\mu g m_{\rm LOS}$ exactly.
\item \textbf{Horizontal balance on tilted surfaces, the Lorentz force in these coordinates, the objective $\Phi$, its weights, the
unknowns} (the pressures below the cut and one height per column), \textbf{the Gauss--Newton iteration, its preconditioners and
step control, and the stopping rule} are those of \code{mhs-poisson-ltau}; only $a$ and its derivative with respect to $\ln P$
differ: $\partial\ln a/\partial\ln P=-\partial\ln\rho/\partial\ln P$, from the equation of state alone.
\end{itemize}
The opacity enters nowhere in the solve. The optical depth is computed afterwards, $d\tau=(\chi/\rho)\,m\,d\ln m$ (as in
\code{column-hse-logm}), and the gauge sets the mean height of $\tau_{500}=1$ to zero.

\section{The cut}
\code{poisson.logm\_cut} $=-0.5$ (default; $m$ in kg\,m$^{-2}$) is the counterpart of $\log\tau=-4$ on the $\tau$ axis: on the Bifrost
snapshot $\log\tau=-4$ has a median $\log_{10}(P/g)=-0.49$, and $\log_{10}m=-0.5$ a median $\log\tau=-3.98$. The two surfaces agree
only in the median, since the opacity varies with $T$. In the snapshot the cut lies about 544\,km above the median $\tau=1$ surface,
near the temperature minimum. With \code{cut\_frame = "vertical"} (default) it is placed at the vertical mass,
$\log_{10}m_{\rm LOS}=$ \code{logm\_cut} $-\log_{10}\mu$ on the ray. Above the cut nothing is solved: $P=gm$ and the levels keep the
hydrostatic heights of the starting state, so the chromosphere is not force-balanced.

\section{The pressure level}
On the $\tau$ axis the pressure level comes from the opacity through the top closure. On the mass axis it is fixed by the axis
itself: at the cut level $P=gm$, which the solve keeps. No reference pressure and no opacity are needed; in exchange the result is
only as good as the input $m$.

\section{What the result means}
On the Bifrost snapshot with $m$ from its own $\rho$ and $P$ (which an inversion does not have), $252^2$ columns at $\mu=1$: the heights
of $\tau=1$ have an rms error of 19.5\,km against the snapshot (column HSE on the same axis: 56.6\,km); the magnetic depression is
$-73.2$\,km against the snapshot's $-76.7$\,km; the run converged in 6 steps. On the $504^2$ cube, cuts from $-2$ to $+0.5$ give almost
the same $z(\tau=1)$ (rms 19.2--19.6\,km), and at $-3$, where 89\,\% of the cube has $\beta<1$, the solve neither converges nor improves.

\section{Assumptions}
Those of \code{mhs-poisson-ltau} (static gas, $T$ and $\B$ fixed on their levels, best compromise rather than equilibrium, rigid
atmosphere above the cut, periodic box, fixed weights, LTE, no condition on $\nabla\cdot\B$), with ``the input $\tau$ scale is
consistent with the table's opacity'' replaced by: \textbf{the input mass axis is correct} (a cgs axis not marked
\code{g cm-2} scales every pressure by 10).

\adjointpoissontrue
% Shared section of the physics documents of the six fixed-B modes: derivatives of the result (the adjoint).
% \input by each mode document after \adjointcolumntrue (column modes) or \adjointpoissontrue (Poisson modes).
\section{Derivatives of the result: the adjoint}\label{sec:adjoint}
An inversion that fits observed spectra changes $T$ (and $\B$) and asks how a misfit $\ell$ that depends on the atmosphere,
through $P$, $\rho$, $n_e$, $\chi_{500}$, $\tau_{500}$ and the heights $z$, changes with them. It needs the gradient
$\partial\ell/\partial\ln T$ at every grid point, often $10^6$ numbers. Perturbing each input in turn and solving again would
cost one or two solves per number. The \textbf{adjoint} gives all of them for the cost of about one solve.

\subsection{The idea}
Write the solve as a map from the inputs $x=(\ln T,\B)$ to the outputs $y$, and the misfit as $\ell(y)$. By the chain rule
\begin{equation}
\frac{\partial\ell}{\partial x}=\Big(\frac{\partial y}{\partial x}\Big)^{\!\top}\bar y,\qquad \bar y\equiv\frac{\partial\ell}{\partial y}.
\label{eq:adj}
\end{equation}
$\bar y$ (the \emph{cotangent} of the outputs) is given by the caller. The Jacobian $\partial y/\partial x$ is never formed:
the transposed product is evaluated by running the computation backwards, each step passing the cotangent of its outputs
to its inputs. This is reverse-mode differentiation, and its cost does not depend on how many inputs there are.

\subsection{What is differentiated}
\ifadjointcolumn
The column integration is a fixed sequence of arithmetic steps (substeps, table look-ups, the top closure, the alignment and
the gauge), so its adjoint is the derivative of exactly that computation: the \textbf{discrete adjoint}. Its discrete choices
are recorded during the solve and held fixed: the number of substeps of each interval, the table cell of each look-up, the
branch of the $\tau$ top closure, the alignment level and its median column, and the interpolation brackets. Away from the
points where one of these choices switches, the map is smooth and this is its exact derivative; at a switch no derivative
exists. Where a step length depends on the state (the Heun substeps of the $\tau$ and mass axes), the length is
differentiated too. An implicitly defined quantity (the hydrostatic top closure) is differentiated by the implicit function
theorem. The equation-of-state derivatives are those of the table interpolant (its bicubic Hermite polynomials), not of the model behind the table.
\fi
\ifadjointpoisson
The Poisson solve is iterative: its result is the state $v^*$ at which the gradient of $\Phi$ vanishes,
$G(v^*,x)\equiv\partial\Phi/\partial v=0$. Differentiating this condition (the \textbf{implicit function theorem}) gives
\begin{equation}
\frac{\partial v^*}{\partial x}=-\mathsf H^{-1}\,\frac{\partial G}{\partial x},\qquad
\mathsf H=\frac{\partial^2\Phi}{\partial v^2},\label{eq:ift}
\end{equation}
so the adjoint solves one linear system with the Hessian $\mathsf H$, $\mathsf H\lambda=\bar v$, and returns
$-(\partial G/\partial x)^{\top}\lambda$, plus the contributions of everything else that depends on $x$: the starting
column-hydrostatic state (which sets the weights of $\Phi$ through their layer medians), the part above the cut and the
gauge. The iterations themselves (damping, step scaling, preconditioners, warm starts) only steer the solver toward $v^*$
and are not differentiated. Two points matter:
\begin{itemize}
\item \textbf{The full Hessian.} With $\B$ fixed the residual does not vanish at the solution (the Lorentz force is in general
not a gradient), so the Gauss--Newton matrix $\mathsf J^{\top}\mathsf W\mathsf J$ of the solver is not the Hessian: the second
derivatives of the residuals contribute too. The adjoint uses the full $\mathsf H$, applied to vectors by differences of the
exact gradient, and solves for $\lambda$ by preconditioned conjugate gradients.
\item \textbf{Directions $\Phi$ does not see.} Some changes of the unknowns leave $\Phi$ unchanged and $\mathsf H$ singular along
them. They keep their starting values, so the derivative passes through to the start along them: on the height grid the mean
$\ln P$ of the anchor plane; on the $\tau$ and mass axes a uniform shift of all heights (fixed by the gauge) and the shortest
horizontal waves of the centred differences, which the differences cannot see.
\end{itemize}
Equation~\eqref{eq:ift} holds only at a stationary point. The Gauss--Newton iteration damps every step and measures its own
progress with that damping, so a slowly converging direction can stay far from stationary while the iteration reports
little left to gain. The option \code{polish} therefore adds up to $n$ \textbf{Newton steps with the exact Hessian} after the
solve, which bring the returned state to $v^*$ (on a small test cube four steps lower the undamped Newton decrement to
$10^{-25}$). The adjoint reports how stationary the returned state was: \code{:stationary} when one undamped Newton step
would lower $\Phi$ by at most a relative $10^{-6}$, \code{:not\_stationary} otherwise, in which case the result is the
derivative of the point the solver was approaching.
\fi

\subsection{How exact it is}
The adjoint is checked against finite differences of the solve itself: $\langle\bar y,\,[y(x+h\,\delta x)-y(x-h\,\delta x)]/2h\rangle$
against $\langle\partial\ell/\partial x,\,\delta x\rangle$ for random $\delta x$ and $\bar y$, on 1, 2 and 3 MPI ranks and every
viewing geometry.
\ifadjointcolumn
The two agree to $3\times10^{-10}$--$3\times10^{-8}$ (relative) on all three axes.
\fi
\ifadjointpoisson
With \code{polish} they agree to $10^{-9}$--$10^{-7}$ (relative).
\fi
A solve with a given electron density (no hydrostatic integration) is forward only and has no adjoint.


\section{Summary of the algorithm}
\begin{enumerate}
\item Read $\log_{10}m$, $x$, $y$, $T$, $\B$. Column HSE on the mass axis (\code{column-hse-logm}) gives the start: $P=gm$, heights,
alignment, gauge.
\item Cut at $\log_{10}m=-0.5$ (vertical). Unknowns: $\ln P$ below the cut and the cut-level height of each column.
\item Minimise $\Phi$ with the Gauss--Newton iteration of \code{mhs-poisson-ltau} ($a=\mu m/\rho$); stop on the qualified
Gauss--Newton decrement, after 25 steps, or when no step is accepted.
\item $\tau_{500}$ from $\kappa\,dm$; gauge the mean $z(\tau{=}1)$ to zero; check the fields; write.
\end{enumerate}
\end{document}
